\documentclass[11pt]{article}
\usepackage{act-sci}
\settitle{The ACT Science Section: Format \& Strategy}
\setsubtitle{Part 1 of 6 \textperiodcentered{} Read this first}
\settarget{$\geq$ 13 correct}
\settime{15 min}

\begin{document}
\maketitleblock

\begin{abstract}
On the enhanced ACT, Science is an optional fourth section: 40 multiple-choice questions in 40 minutes, built on six or seven passages of graphs, tables, experiment summaries, and one set of conflicting viewpoints. It is a reasoning test that happens to use science as its material. Most questions can be answered from the page, and the main threats are reading too slowly, misreading a figure, and losing time on one hard question. This part covers how the section is built and scored, a pacing plan with checkpoints, a passage-reading method, a map from question wording to the skill being tested, and the answer-choice traps ACT repeats. It ends with a classification drill, a pacing drill, and a short passage to try the method on. The other five parts are listed in Section~\ref{sec:plan}.
\end{abstract}

\scoringline

\tableofcontents
\clearpage

%% =====================================================================
\section{Rule OV-1: Know exactly what you are walking into}

\begin{rulebox}{OV-1}{The section at a glance}
\renewcommand{\arraystretch}{1.15}
\begin{tabularx}{\linewidth}{@{}l X@{}}
Questions / time & 40 questions, 40 minutes ($\approx$ 60 s per question). Optional; taken after English, Math, and Reading.\\
Passages & Usually 6--7, each with 5--7 questions. Biology, chemistry, physics, Earth and space science, and (new) engineering design.\\
Answer choices & Four per question, lettered A--D on odd-numbered questions and F--J on even-numbered ones.\\
Scoring & Scaled 1--36, reported separately. Science is \emph{not} part of the composite but is averaged with Math into the STEM score. Some questions are unscored field-test items, and you cannot tell which.\\
Guessing & No penalty for wrong answers. Never leave a question blank.\\
Outside knowledge & A small number of questions assume basic science facts (Part 5). Everything else is on the page.\\
\end{tabularx}
\end{rulebox}

\subsection*{Three passage formats}
\small
\begin{tabularx}{\linewidth}{@{}l X l@{}}
\toprule
Format & What it looks like & Skills it favours\\\midrule
Data Representation & A short intro and 2--4 graphs or tables. Little text. & IOD (Part 2)\\
Research Summaries & Two or three experiments with procedures and results. & SI (Part 3), IOD, EMI\\
Conflicting Viewpoints & 2--4 scientists or students explain one phenomenon differently. Mostly text. One per test. & EMI (Part 4)\\
\bottomrule
\end{tabularx}
\normalsize

\subsection*{Three reporting categories}
\small
\begin{tabularx}{\linewidth}{@{}l c c X@{}}
\toprule
Category & Share & $\approx$ Questions & Core question\\\midrule
Interpretation of Data (IOD) & 36--47\% & 15--19 & What do the figures say?\\
Scientific Investigation (SI) & 17--31\% & 7--12 & How was the experiment designed and done?\\
Evaluation of Models, Inferences \& Results (EMI) & 22--36\% & 9--14 & What do the results mean, and who is right?\\
\bottomrule
\end{tabularx}
\normalsize

%% =====================================================================
\section{Rule OV-2: Pace by passage, not by question}

\begin{rulebox}{OV-2}{About five and a half minutes per passage}
Forty minutes over seven passages is about 5\,min 40\,s each. Check the clock only at passage boundaries, against these checkpoints:
\begin{center}
\renewcommand{\arraystretch}{1.1}
\begin{tabular}{@{}l c c c c c c c@{}}
\toprule
Passages finished & 1 & 2 & 3 & 4 & 5 & 6 & 7\\\midrule
Elapsed time (7 passages) & 5:40 & 11:20 & 17:00 & 22:50 & 28:30 & 34:10 & 40:00\\
\bottomrule
\end{tabular}
\end{center}
If you are more than a minute behind, take one guess on the hardest remaining question in the next passage instead of working it out. In the last minute, fill in every blank.
\end{rulebox}

\trap{The biggest time loss is not a slow passage but one stubborn question. If 90 seconds pass without an answer, guess, mark it, and move on. One question is worth the same as the three easy ones you would lose to it.}

%% =====================================================================
\section{Rule OV-3: Choose your order}

\begin{rulebox}{OV-3}{Easy points first, Conflicting Viewpoints last}
\begin{enumerate}
\item Skim the whole section for 20 seconds. Identify the Conflicting Viewpoints passage (a block of paragraphs headed \emph{Scientist 1}, \emph{Student 2}, \emph{Hypothesis 1}\ldots).
\item Work the Data Representation passages first. They are short and mostly IOD.
\item Then the Research Summaries.
\item Do Conflicting Viewpoints last, unless reading is your strength. It takes the most reading time but rewards a careful read.
\end{enumerate}
\end{rulebox}

%% =====================================================================
\section{Rule OV-4: Skim the passage, then let questions send you back}

\begin{rulebox}{OV-4}{The 40-second skim}
\textbf{Data and research passages:}
\begin{enumerate}
\item Read the first paragraph (the topic and the purpose).
\item For each experiment, read only enough to note what was \emph{varied} and what was \emph{measured} (SI-4).
\item Read the \emph{frame} of each figure: title, axes, units, legend (ID-1). Do not read the data yet.
\item Go to the questions. Each question tells you where to look (``According to Figure 2\ldots'').
\end{enumerate}
\textbf{Conflicting Viewpoints:} read the intro and each viewpoint \emph{fully}, writing a one-line summary after each (EM-3). Then answer.
\end{rulebox}

%% =====================================================================
\section{Rule OV-5: Decode the question stem}

\begin{rulebox}{OV-5}{The wording tells you the skill}
\small
\renewcommand{\arraystretch}{1.12}
\begin{tabularx}{\linewidth}{@{}>{\raggedright\arraybackslash}X l l@{}}
Stem wording & Category & Rule\\\midrule
``According to Figure/Table $n$, \ldots{} was closest to'' & IOD & ID-1\\
``As $x$ increased, $y$\ldots'' / ``Which best describes the relationship'' & IOD & ID-2\\
``If an additional trial were done at\ldots'' / ``would most likely be'' & IOD / EMI & ID-3, EM-2\\
``Based on Figures 1 and 2\ldots'' & IOD & ID-4\\
``Which was the independent / dependent / controlled variable?'' & SI & SI-1\\
``The purpose of the control / of rinsing / of repeating\ldots'' & SI & SI-2, SI-3\\
``Experiments 1 and 2 differed in that\ldots'' & SI & SI-4\\
``Which procedure would best test\ldots'' & SI & SI-5\\
``Which method was more precise / accurate'' & SI & SI-6\\
``Which prototype meets the criteria\ldots'' & SI & SI-7\\
``Do the results support the hypothesis\ldots?'' & EMI & EM-1\\
``According to Scientist 2\ldots'' & EMI & EM-3\\
``Scientists 1 and 2 would both agree\ldots'' & EMI & EM-4\\
``Which finding would strengthen / weaken\ldots'' & EMI & EM-5\\
``Which conclusion is best supported\ldots'' & EMI & EM-6\\
\end{tabularx}
\end{rulebox}

%% =====================================================================
\section{Rule OV-6: Eliminate the six standard traps}

\begin{rulebox}{OV-6}{Wrong answers are built from a short list}
\begin{enumerate}
\item \textbf{Right number, wrong place:} a value from the neighbouring row, the other series, or the other figure.
\item \textbf{Reversed trend:} ``increases'' where the data decrease, or an ``only'' choice that misses a peak.
\item \textbf{True but irrelevant:} a correct fact that does not answer the question (EM-1).
\item \textbf{Absolute language:} ``always,'' ``proves,'' ``for all,'' ``never'' (EM-6).
\item \textbf{Unit or scale slip:} \si{\milli\litre} vs.\ \si{\litre}, linear reading of a log axis (ID-6).
\item \textbf{Outside knowledge over the passage:} a choice that is true in real life but contradicts the passage or the viewpoint being asked about.
\end{enumerate}
When you guess, eliminate first: removing one or two traps raises the odds of a correct guess from 25\% to 33--50\%.
\end{rulebox}

%% =====================================================================
\section{How the series fits together}\label{sec:plan}

\small
\begin{tabularx}{\linewidth}{@{}l l X l@{}}
\toprule
Part & Title & Contents & Questions\\\midrule
1 & Format \& Strategy & This document. & 15\\
2 & Interpretation of Data & Rules ID-1 to ID-7: reading values, trends, interpolation, linking figures, math, units, graph matching. & 26\\
3 & Scientific Investigation & Rules SI-1 to SI-7: variables, controls, procedures, design, error, engineering. & 23\\
4 & Evaluation of Models & Rules EM-1 to EM-6: hypotheses, prediction, conflicting viewpoints, strengthen/weaken, scope. & 20\\
5 & Science Knowledge Primer & The outside facts ACT assumes, by subject. & 24\\
6 & Full-Length Practice Test & 40 questions, 40 minutes, 7 passages, blueprint-matched. & 40\\
\bottomrule
\end{tabularx}
\normalsize

\subsection*{Four-week plan}
\small
\begin{tabularx}{\linewidth}{@{}l X@{}}
\toprule
Week & Work\\\midrule
1 & Part 1, then Part 2 untimed. Redo every missed question and write the rule it tested.\\
2 & Parts 3 and 4 untimed. Redo the Part 2 drills timed (60 s per question).\\
3 & Part 5 (one subject per day). Redo Parts 3--4 timed.\\
4 & Part 6 under real conditions (40 min, one sitting). Review using the Rule column in the key; return to the guide for every rule you missed twice.\\
\bottomrule
\end{tabularx}
\normalsize

%% =====================================================================
\section{Drill A: name the skill}

\tip{For each stem, choose the reporting category it belongs to. This trains you to know which rule to use before you read the choices.}

\begin{q}{1}{OV-5}{Reading a value from a figure is IOD (ID-1).}
``According to Figure 2, at \SI{40}{\celsius} the solubility of KNO$_3$ is closest to:''
\choicesgrid{Interpretation of Data}{Scientific Investigation}{Evaluation of Models}{Outside knowledge only}
\end{q}

\begin{q}{2}{OV-5}{Identifying a controlled variable is about experimental design (SI-1).}
``Which of the following was a controlled variable in Experiment 2?''
\choicesgrid{Interpretation of Data}{Scientific Investigation}{Evaluation of Models}{Outside knowledge only}
\end{q}

\begin{q}{3}{OV-5}{Agreement between viewpoints is EMI (EM-4).}
``Scientists 1 and 2 would most likely agree that the Moon:''
\choicesgrid{Interpretation of Data}{Scientific Investigation}{Evaluation of Models}{Outside knowledge only}
\end{q}

\begin{q}{2}{OV-5}{The purpose of a procedural step is SI (SI-3).}
``The students rinsed the beaker with distilled water between trials most likely to:''
\choicesgrid{Interpretation of Data}{Scientific Investigation}{Evaluation of Models}{Outside knowledge only}
\end{q}

\begin{q}{3}{OV-5}{Judging whether results support a hypothesis is EMI (EM-1).}
``Do the results of Experiment 1 support the hypothesis that enzyme activity is greatest at \SI{37}{\celsius}?''
\choicesgrid{Interpretation of Data}{Scientific Investigation}{Evaluation of Models}{Outside knowledge only}
\end{q}

\begin{q}{1}{OV-5}{Describing a trend in the data is IOD (ID-2).}
``As altitude increased, air pressure:''
\choicesgrid{Interpretation of Data}{Scientific Investigation}{Evaluation of Models}{Outside knowledge only}
\end{q}

\begin{q}{3}{OV-5}{Deciding which viewpoint a new finding weakens is EMI (EM-5).}
``Suppose researchers found that Moon rocks were never heated above \SI{500}{\celsius}. This finding would weaken the viewpoint of:''
\choicesgrid{Interpretation of Data}{Scientific Investigation}{Evaluation of Models}{Outside knowledge only}
\end{q}

\begin{q}{4}{OV-5}{Nothing in a passage is needed: the formula of table salt (NaCl) is a fact you bring with you (Part 5).}
``Which of the following is the chemical formula of table salt?''
\choicesgrid{Interpretation of Data}{Scientific Investigation}{Evaluation of Models}{Outside knowledge only}
\end{q}

%% =====================================================================
\section{Drill B: pacing decisions}

\begin{q}{3}{OV-2}{The checkpoint for 3 passages is about 17:00. At 18:00 you are about one minute behind: recoverable with a slightly faster next passage.}
On a 7-passage section, you finish the third passage when 18 minutes have elapsed. You are:
\choices{well ahead of pace.}{exactly on pace.}{about one minute behind; work slightly faster on the next passage.}{so far behind that you should skip the remaining passages.}
\end{q}

\begin{q}{4}{OV-1}{There is no penalty for wrong answers, so a blank is a guaranteed zero and a guess has at least a 25\% chance.}
With one minute left, six questions are unanswered. The best move is to:
\choices{work the easiest of the six carefully.}{leave them blank to avoid losing points.}{reread the last passage's introduction.}{enter an answer for all six.}
\end{q}

\begin{q}{2}{OV-3}{Conflicting Viewpoints needs the most reading per question, so most students should do it last.}
For most students, which passage type is best left for last?
\choicesgrid{Data Representation}{Conflicting Viewpoints}{Research Summaries}{Whichever comes first}
\end{q}

\begin{q}{1}{OV-2}{90 seconds is the cut-off. Guess, mark, move, and return only if time remains.}
You have spent 90 seconds on a question without settling on an answer. You should:
\choices{guess, mark the question, and move on.}{keep working until you are sure.}{skip the rest of the passage.}{reread the whole passage.}
\end{q}

%% =====================================================================
\section{Drill C: try the method}

\tip{Use the 40-second skim: one sentence of intro, then the frame of Table 1, then the questions.}

\passage{Rainfall and river discharge}
Hydrologists recorded the total rainfall and the average discharge (volume of water flowing past a point each second) of a river during six months (Table 1).

\figlabel{Table 1}
\datatable{l c c}{Month & Rainfall (mm) & Discharge (\si{\metre\cubed\per\second})\\\midrule
January & 120 & 340\\ February & 95 & 310\\ March & 80 & 260\\ April & 40 & 180\\ May & 15 & 110\\ June & 5 & 70\\}
\passageend

\begin{q}{2}{ID-1}{Row ``April,'' column ``Discharge'': 180.}
According to Table 1, the river's discharge in April was:
\choicesrow{\SI{110}{\metre\cubed\per\second}}{\SI{180}{\metre\cubed\per\second}}{\SI{260}{\metre\cubed\per\second}}{\SI{340}{\metre\cubed\per\second}}
\end{q}

\begin{q}{1}{ID-2}{Both columns fall every month from January to June: as rainfall decreased, discharge decreased.}
According to Table 1, as monthly rainfall decreased, discharge:
\choices{decreased only.}{increased only.}{decreased and then increased.}{stayed the same.}
\end{q}

\begin{q}{3}{ID-3}{60 mm lies between April (40 mm $\to$ 180) and March (80 mm $\to$ 260); halfway gives about 220.}
In a month with \SI{60}{\milli\metre} of rainfall, the discharge would most likely be closest to:
\choicesrow{\SI{110}{\metre\cubed\per\second}}{\SI{150}{\metre\cubed\per\second}}{\SI{220}{\metre\cubed\per\second}}{\SI{300}{\metre\cubed\per\second}}
\end{q}

\makekey
\end{document}
